% ---- pg-025 （书页 13） ----

的相似(和不同)之处.项目的分组会在这个表象中表现出来.

下面的例子说明三维散布图.

\mvbold{例 1.6}\quad （寻找降维结构）

一位动物学家得到了 $n=25$ 组蜥蜴的数据.其重量(体重)的单位为克,嘴到排泄口的体长(SVL)和后肢跨距(HLS)的单位为毫米.\hyperref[tab:1.3]{表 1.3} 中列出了对应的数据.

\begin{mvtable}
\begin{tabular}{cccc|cccc}
\toprule
蜥蜴 & 体重 & SVL & HLS & 蜥蜴 & 体重 & SVL & HLS \\
\midrule
1 & 5.526 & 59.0 & 113.5 & 14 & 10.067 & 73.0 & 136.5 \\
2 & 10.401 & 75.0 & 142.0 & 15 & 10.091 & 73.0 & 135.5 \\
3 & 9.213 & 69.0 & 124.0 & 16 & 10.888 & 77.0 & 139.0 \\
4 & 8.953 & 67.5 & 125.0 & 17 & 7.610 & 61.5 & 118.0 \\
5 & 7.063 & 62.0 & 129.5 & 18 & 7.733 & 66.5 & 133.5 \\
6 & 6.610 & 62.0 & 123.0 & 19 & 12.015 & 79.5 & 150.0 \\
7 & 11.273 & 74.0 & 140.0 & 20 & 10.049 & 74.0 & 137.0 \\
8 & 2.447 & 47.0 & 97.0 & 21 & 5.149 & 59.5 & 116.0 \\
9 & 15.493 & 86.5 & 162.0 & 22 & 9.158 & 68.0 & 123.0 \\
10 & 9.004 & 69.0 & 126.5 & 23 & 12.132 & 75.0 & 141.0 \\
11 & 8.199 & 70.5 & 136.0 & 24 & 6.978 & 66.5 & 117.0 \\
12 & 6.601 & 64.5 & 116.0 & 25 & 6.890 & 63.0 & 117.0 \\
13 & 7.622 & 67.5 & 135.0 & & & & \\
\bottomrule
\end{tabular}
\phantomsection\label{tab:1.3}\mvtabcap{表 1.3\quad 蜥蜴数据}
\end{mvtable}

{\zihao{6} 资料来源：数据经 Kevin E. Bonine 许可使用.}

尽管有 3 种大小的测量数据,我们感兴趣的是大部分的变化是否都基本限制在二维甚至于一维空间上.

为了回答上面关于降维的问题,在\hyperref[fig:1.6]{图 1.6} 中我们画出了数据的三维散布图.很明显大部分的变化都集中在一条一维的直线上.知道了沿这些点的主轴直线的位置,也就大致知道了这三维数据(重量,SVL 和 HLS)的情况.

可是,如果在一个变量的变化比其他的变化大得多的情况下,这种分析会误导大家.所以,我们首先要将数值标准化,即 $z_{jk} = (x_{jk} - \bar{x})/\sqrt{s_{kk}}$,这样各个变量在散布图中有相同的作用.\hyperref[fig:1.7]{图 1.7} 为标准化后变量的散布图,大部分的变化都能被由这些点决定的直线所对应的单个变量来解释.

\begin{mvfigure}
\mvinc{pg-025}{0.11957,0.72689,0.48374,0.89050}
\phantomsection\label{fig:1.6}\mvcap{图 1.6\quad 表 1.3 的蜥蜴数据的三维散布图}
\end{mvfigure}

\begin{mvfigure}
\mvinc{pg-025}{0.52145,0.72819,0.89459,0.89050}
\phantomsection\label{fig:1.7}\mvcap{图 1.7\quad 标准化后数据的三维散布图\quad$\blacksquare$}
\end{mvfigure}

% UNCERTAIN: pg-025 \hyperref[fig:1.7]{图 1.7} 图题行右端实心方块 ■ 为原书"节/例结束"符号，已用 $\blacksquare$ 还原（参 pg-017 先例）。
